\section{Related Works}
\label{sec:relwork}
% \hh{In this study, we aim to extract deep support vectors from a deep network by utilizing the equivalence between the deep learning model and SVM. This section will explain the foundational papers on which our idea is based, and how our work can be interpreted.}

% \jh{우리는 Deep learning model과 SVM의 등가성을 이용해 딥 러닝에서 deep support vector를 추출하고자 한다. Self-containedness를 위해서 본 섹션에서는 우리의 아이디어가 어떤 논문에 기반을 두고 있는지, 또한 우리의 논문이 어떤 식으로 해석될 수 있는지를 설명하도록 한다. }ㅐ


% [1]Deep Total Variation Support Vector Networks, 10.1109/ICCVW.2019.00365
% [2]When Ensemble Learning Meets Deep Learning: a New Deep Support Vector Machine for Classification, 10.1016/j.knosys.2016.05.055 [3]An Ensemble of Deep Support Vector Machines for Image Categorization, 10.1109/SoCPaR.2009.67
% [4]Deep Support Vector Classification and Regression, 10.1007/978-3-030-19651-6_4
% [5]A Complete Deep Support Vector Data Description for One Class Learning, 10.1109/ACCESS.2023.3325734
% [6]Deep support vector machine for hyperspectral image classification, 10.1016/j.patcog.2020.107298
% [7]Contrastive deep support vector data description, 10.1016/j.patcog.2023.109820
% [8]Deep neural mapping support vector machines, 10.1016/j.neunet.2017.05.010
% [9]Deep Learning using Linear Support Vector Machines, https://doi.org/10.48550/arXiv.1306.0239 (https://doi.org/10.48550/arXiv.1306.0239) [10]Totally Deep Support Vector Machines, https://doi.org/10.48550/arXiv.1912.05864 (https://doi.org/10.48550/arXiv.1912.05864)



\subsection{SVM in Deep Learning}
Numerous studies have endeavored to establish a connection between deep learning and conventional SVMs. 
\jh{In the theoretical side}, \cite{svm:theory} demonstrated that, in an overparameterized setting with linearly separable data, the learning dynamics of a model possessing a logistic tail are equivalent to those of a support vector machine, with the model’s normalized weights converging towards a finite value. Following this, \cite{svmhomo} extended this equivalence to feedforward networks. \jh{\nj{This} line of research relies on strong assumptions such as full-batch training \nj{and} non-residual architecture \nj{without data augmentation}.}
% \jh{There exists a line of research on integrating SVM architecture into deep learning, which is called DeepSVM to \nj{benefit from its desirable} properties.} \cite{svmpractical, svmpractical2, SVMex1,SVMex2,SVMEX3}. \jh{DeepSVM integrated SVMs as feature extractor\nj{s,} extracting meaningful feature\nj{s in a} human-crafted manner.} 
% Despite these advancements, there is still a lack of research that directly connects support vectors through this theoretical lens of equivalence. In this study, we address this theoretical equivalence of support vectors %head-on 
% by introducing the DeepKKT condition, a KKT condition tailored for deep learning, thus applying this equivalence in a practical context.
% \jh{We \nj{show} that reconstructing support vectors in deep models is \nj{indeed} possible \nj{and obtaining high-quality} support vectors is feasible under much smoother conditions.}
\jh{There also exists a body of work on integrating SVM principles into deep learning, often referred to as DeepSVM, aiming to leverage \nj{SVM’s desirable} properties} \cite{svmpractical, svmpractical2, SVMex1,SVMex2,SVMEX3}. \jh{DeepSVM integrates SVM components, specifically using them as feature extractors to derive meaningful, human-crafted features.}

In contrast, our work does not modify or incorporate SVM architectures. Instead, we focus on identifying support vectors directly within deep learning models, thereby bridging the gap between deep learning and support vector machines in a more fundamental manner. Despite these advancements, there remains a lack of research that directly connects support vectors through a theoretical lens of equivalence. In this study, we address this gap by introducing the DeepKKT condition, a KKT condition tailored for deep learning, allowing us to apply the concept of support vectors in a practical deep learning context.

\jh{We \nj{show} that reconstructing support vectors in deep models is indeed feasible, and obtaining high-quality support vectors is achievable under much less restrictive conditions compared to prior work.}



\subsection{\jh{Model Inversion Through the Lens of Maximum Margian}}

There is a line of research utilizing the stationarity condition, \nj{a} part of the KKT condition, for model inversion. \cite{svm:marginreconstruct} firstly exploited the KKT condition for model generation, adopting SVM-like architectures. They \jh{normally} conducted experiments with binary classification, a 2-layer MLP, and full-batch gradient descent. \cite{svm:reconstructmulti} extended these experiments to multi-label classification by adapting their existing architecture to a multi-class SVM structure. To ensure the generated samples lie on the data manifold, 
% \jh{they were initialized with the dataset's mean,}
\jh{they initialized with the dataset's mean,}
implying the adoption of some prior knowledge of the data. Similarly, \cite{margin:genfclassifer} generated images through the stationarity condition, also adopting the mean-initialization and conducting experiments on the \nj{CIFAR10 \cite{dataset:cifar10},} MNIST \cite{data:mnist} and \jh{downsampled} CelebA~\cite{dataset:celebA} datasets. 
\nj{Their work generally focused on low-dimensional, labeled datasets with a small number of classes such as CIFAR10 and MNIST, consistent with the traditional SVM setting of binary-labeled, low-dimensional datasets. In contrast, we extended our experiments to high-dimensional datasets with many classes, an area traditionally dominated by deep learning. Specifically, we conducted experiments on ImageNet~\cite{imagenet} using a pretrained ResNet50 model following the settings described in the original paper \cite{resnet}.}

%\jh{Their work generally concentrated on low-dimensional, small-number of classes labeled datasets such as CIFAR10 %\cite{dataset:cifar10} 
%and MNIST, %\cite{data:mnist}, 
%which aligns with the traditional SVM setting -- binary-labeled, low-dimensional datasets. However, we reconstructed the data in high-dimensional, large-number of class\nj{es} settings, which is an area traditionally explored by deep learning. We conducted experiments on the most \nj{popular} settings in the deep learning realm, specifically on ImageNet \cite{imagenet}, using a pretrained ResNet50 following the original paper's settings \cite{resnet}.}

\nj{Furthermore, previous works have concentrated on \textbf{reconstructing the training} dataset. In contrast, similar to generative models, our work focuses on \textbf{generating unseen} data from noise using a classification model. Additionally, we emphasize the original meaning of \hh{`support vectors'.} Unlike other approaches, our Deep Support Vectors (DSVs) adhere to the traditional role of support vectors: they explain the decision criteria, and a small number of DSVs can effectively reconstruct the model.}

%\jh{Furthermore, previous works \nj{have} focused on \textbf{reconstructing} the `training' dataset. In contrast, similar to generative models, our work \nj{focuses} on \textbf{generating} `unseen' datasets from noise using a classification model.}
%\jh{Additionally, we \nj{emphasize} the original meaning of `support vectors.' Unlike other works, our DSVs (Deep Support Vectors) adhere to the role of support vectors\nj{; DSVs} can explain the \nj{decision} criterion \nj{and a small number of DSVs can reconstruct the model. }} % with a small number of DSVs.}

% \jh{Their works generally concentrated on low-dimensional, small-number of class labeled dataset such as CIFAR10\cite{dataset:cifar10} and MNIST\cite{data:mnist}, which is the setting of traditional SVM --binary-labeled, low-dimensional--. Yet, we reconstructed the data alongside high-dimensional-large-number of classed settings- which is an original area of deep learning. We conducted experiments on most frequent settings in deep learning realm. We conducted on ImageNet \cite{imagenet}, with predrained ResNet50 following original paper's settings\cite{resnet}.}
% \jh{Furthermore, previous works focused on \textbf{reconstructing} `training' dataset. However, similar to generative model our work focused on \textbf{generating} `unseen' dataset from noise with classification model.}
% \jh{Also, we focused on the original meaning `support vector'. Contrary to other works, our DSVs follow the role of support vectors. It can explain the criterion, reconsturct the model within small number of DSVs}
% Contrary to previous studies, our research extensively employs the KKT conditions, not just focusing on the stationarity condition. While previous works adopted SVM-like structures to directly utilize the stationarity condition,
%our experiments are based on a more general setting. 
% derived DSVs from \jh{totally} universal deep learning settings including SGD (Stochastic Gradient Descent), mini-batch processing, multi-class classification, and data augmentation, indicating that DSVs can be obtained in any trained model settings, not just within SVM-like architectures. 
% Furthermore, their limited number aligns with the core principle of support vectors. These distinctions underscore the unique contributions of our work. This was possible because we reinterpreted the KKT condition used in machine learning to fit in deep learning settings. 
% Unlike previous \nj{works, not only did we utilize the stationarity condition but we made use of other conditions except complementary slackness} to ensure that DSVs satisfy them. 
% In contrast to prior studies, our approach goes beyond leveraging the stationarity condition alone; we incorporate additional conditions to ensure DSVs meet them.
% Additionally, we extracted DSVs from random noise, an approach that does not rely on any prior knowledge of the training dataset. Most significantly, we have integrated the concept of conventional support vectors into the realm of deep learning and have formulated the definition of DSVs, which are assumed to play a similar role as conventional support vectors in SVM.

% \hh{Contrary to previous studies, our research extensively employs the KKT conditions rather than solely focusing on the stationary condition, Also, they adopted SVM-like structure to directly utilize stationarity condition. However, 
% our experiments are based on most general setting, DSVs adhere} to the foundational deep learning settings of SGD, mini-batch, multi-class classification and augmentation while their limited number aligns with the core principle of support vectors. These distinctions highlight the unique contributions of our work compared to the previous study. \jh{이것은 우리가 machine learning에서 사용되던 KKT condition을 deep learning settings에 맞게 새롭게 재해석했기 때문에 가능했으며, 우리는 기존의 paper와는 달리 stationarity condition이 아닌 다른 모든 조건들을 활용하여 DSVs가 그 조건을 만족하도록 하였다.}
% Additionally, we extracted DSVs \nj{from} random noise, an approach that does not rely on any prior knowledge of the training dataset.  Most significantly, we have integrated the concept of conventional support vectors into the domain of deep learning and have formulated the definition of DSVs, which \nj{take} a similar role.}


% \jh{여기서 기존의 maximum-margin based paper과 연관지어서 설명하기.}


% \subsection{\jh{Model inversion}}

% \jh{여기서 Dataset Distillation이나,  adversarial attack을 추가하자. }

% \jh{\subsection{Generation Model}
% Sample을 찾는다는 측면에서 유사하다. Boundary를 찾는다. 이런 말들을 넣자.}
% \jh{위 섹션에서 설명했던 것처럼, 딥 러닝에서 등가성이 있는데, 그리고 이것이 KKT condition을 만족한다는 뜻, 관련 논문 reference. 첫 논문은 overparameterized setting에서 linearly separable data에서 logistic tail을 가지고 있을 때, model의 learning dynamics가 support vector machine임과 등가임을 증명하였고, 이 다음 어떤 페이퍼가 feedforward network까지도 SVM과 등가임을 보여주었다. 이런식으로 이론적으로 equivalence를 보이는 페이퍼가 있었고, 반대로 SVM의 좋은 성질을 직접적으로 활용하기 위해서 SVM을 직접으로 아키텍쳐에 접목한 페이퍼도 있다. 그럼에도 불구하고, 현재 이론적 등가성을 이용하여 support vector를 직접적으로 연구한 페이퍼는 없는 실정이다. 이 중 등가성에서 KKT condition을 이용한 페이퍼가 기존에 하나 존재하였지만, 이들의 목적은 KKT condition을 만족하여서 support vector를 만드는 것이 아닌 data manifold 안에 있는 이미지를 generation하는 연구였다. 또한 이들은 일반적인 딥러닝 세팅이 아니라 SVM의 성질을 곧바로 가져오기 위해서 SVM condition과 유사한 binary classifcation, 2 Layer MLP, full-batch gradient descent를 사용하였다. 이후 페이퍼(워크샵)이 multi-class를 사용하였지만, 역시 유사한 문제는 견지하고 있었으며, 이들은 manifold 상에서 이미지를 생성하기 위해서 dataset의 mean에서 시작하였고 수천개의 image를 generate 하는것을 메인 방법론으로 삼았다. 우리는 반대로 support vector를 직접 generate하는 것을 목표로 하며, noise에서 시작하고 manifold상에 있다고 가정하는 것이 아니라, KKT condition을 만족하는 것만으로 manifold상에 있을 수 없음을 지적하고 manifold 상의 거리를 재는 metric을 주장한다.}
% \jh{여기서는 SVM이 어떤 건지를 recap 하는 것, 그리고 이론적 연구, 그리고 이걸 사용하려는 연구는 있었지만, 모델을 직접적으로 svm으로 만들려고 하는 연구는 없었음을 밝힘 기존의 kkt 쓴거는 main goal이 우리의 목표와는 아예 정반대였음을 지적할 것.}
% \lipsum[1-1]
\subsection{Dataset Distillation}
Dataset distillation \cite{distillation_trajectory,distillation:grad1,distillation:nongrad1} fundamentally aims to reduce the size of the original dataset while maintaining model performance. 
\jh{The achievement also involves addressing privacy concerns and alleviating communication issues between \nj{the} server \nj{and client} nodes. }
% \jh{Although \nj{it} does not explicitly clarify the link with SVM, distilled datasets share key features with support vectors in that they enable model reconstruction with minimal data. }
% \jh{The problem is typically addressed under the following conditions: 1) Access to the entire dataset for gradient matching, 2) Possession of snapshots from all stages of the model's training phase, which is impractical setting concerning practical condition\cite{practicalDD}.}
% \jh{In SVM, one can reconstruct original model with support vectors, furthermore, this reconsturction does not pose any condition mentioned above; which is much practical compared to previous dataset distillation settings. DSVs can also serve similar role to support vectors within practical condition. DeepKKT condition does not require none of above conditions.}
The \nj{dataset distillation} problem is typically addressed under the following conditions: 1) Access to the entire dataset for gradient matching, 2) Possession of snapshots from all stages of the model's training phase, which \nj{are impractical settings} \nj{for practical usage} \cite{practicalDD}.  \jh{Furthermore, these algorithms typically }\hh{require Hessian computation, which imposes a heavy computational burden.}

\jh{In SVM, the model can be reconstructed using support vectors. This reconstruction is more practical compared to previous dataset distillation methods, as it does not require any of the restrictive conditions. }
\jh{Likewise, because Deep Support Vectors (DSVs) also do not require these conditions and \nj{are Hessian-free}, they can play the role of distillation under practical conditions.}

% \nj{In SVM, the original model can be reconstructed using support vectors. This reconstruction is more practical compared to previous dataset distillation methods, as it does not require any of the restrictive conditions mentioned earlier. }
% \jh{Likewise, because Deep Support Vectors (DSVs) also do not require these conditions and do not require the Hessian, they can play the role of support vectors under practical conditions.}

%Likewise, because  Deep Support Vectors (DSVs) also do not require these conditions \jh{and does not require the Hessian}, they can play the role of support vectors under practical conditions.} 
%\jh{In SVM, one can reconstruct the original model with support vectors. Furthermore, this reconstruction does not pose any of the conditions mentioned above, making it much more practical compared to previous dataset distillation settings. DSVs can also serve a similar role \nj{as} support vectors under practical conditions. The DeepKKT condition does not require any of the above conditions.}

% These conditions present greater complexity and constraints relative to \nj{our} DSV extraction, posing practical implementation challenges in real-world situations. In contrast, our algorithm does not require all three of these conditions and operates \nj{based just on} a single trained model.
% \jh{Dataset distillation은 원본 데이터셋의 크기를 줄이면서도 모델의 성능을 유지하는 것을 기본 골자이며, 해당 문제의 성취를 위해서 privacy 문제를 해결하고 서버 노드 간 communication 문제를 해소하는 것을 목표로 한다. 비록 해당 문제가 SVM과의 연관성을 푸는 것은 아니지만, 결국 모델을 재구축할 수 있는 적은 데이터셋을 푼다는 점에서, Distilled dataset은 Support vector가 가져야 하는 성질을 가지고 있다. 이들은 주로 다음과 같은 조건 속에서 문제를 푸는데, 1. Gradient Matching을 해야 하기 때문에 전체 데이터셋을 볼 수 있고, 2. 모델의 training 단계 전체의 snapshot을 다 가지고 있다. 3. 서로 다른 initialization 한 모델 여러개를 가지고 있다. 이러한 조건들은 자연스럽게 DSV가 추출되어야 하는 조건을 위배하며, 우리의 알고리즘은 위 세 조건 전부 필요로 하지 않으며, 단지 학습된 모델 하나만 가지고 있는것을 전제한다.}

% \subsection{Responsible AI}

% Responsive Artificial Intelligence (RAI) is an approach that emphasizes ethical considerations in AI development \cite{RAI:definition}. This field focuses on addressing the issues that arise from the models' indifference, as they mainly consider numerical criteria. The `black-box' nature of deep learning models presents potential risks. For instance, data is often centered around Western culture, leading to possible racial biases in models \cite{RAI:western}.
% However, the focus is primarily on large models such as Large Language Models (LLMs) and diffusion models~\cite{RAI:LLM, RAI:diffusion} with fewer concerns directed towards classification-variant models, which make up a significant portion of deep learning models. While Explainable AI (XAI) \cite{montavon2017explaining} addresses this aspect, it often lacks global explainability.

% Deep Support Vectors (DSVs) naturally meet the conditions that RAI concerns. DSVs enable the visualization of models' decision criteria at a low cost. Unlike conventional XAI methodologies, which primarily rely on local explanations \cite{RAI:Gradcam, XAIrescent} (explanations within a given data set), DSVs do not depend on the training dataset. This independence implies enhanced safety and privacy, providing explainability to the displayed model. With DSVs, we can assess which models embody fairness based on qualitative criteria.

% \subsection{Explainable AI (XAI)}

% \hh{Explainable Artificial Intelligence (XAI) is a field developed to address the inherently complex and non-transparent nature of deep learning models. This field concentrates on interpreting and elucidating the underlying principles of the results generated by deep learning models, especially in the analysis of data such as images. This is known as local explanation where explanations are conditional on the data. In contrast, global explanations describe the overall decision-making process of a model, which does not rely on specific data. Until now, this area has mostly been theoretical with no scholarly work turning it into a practical reality. Support vectors are the data points closest to the decision boundaries of \nj{a model} and they implicitly represent the decision-making processes of \nj{the model} across the overall data manifold. This study shows that support vectors initially designed for linear spaces can be effectively used in nonlinear deep\nj{-}learning models to provide strong and efficient global explanations.}
% % \jh{XAI는 딥 러닝 모델이 가지고 있는 blackbox 적인 특성을 해소하기 위해서 연구되고 있는 분야로, 주로 이미지가 주어졌을 때, 어떤 예측을 기반으로 이루어졌는지를 설명한다. 이를 local explanation이라고 하며, 이들의 설명은 data에 conditional하게 이루어진다. 반면에 data에 conditional하지 않게 모델 전반의 의사 결정 절차를 설명하는 것을 global explanation이라고 하며, 이때까지는 해당 논의가 이론상으로만 존재하고, 실제로 이를 가능케 한 논문이 존재하지 않았다. 본 논문에서는, Support Vector가 가진 본래적인 정의를 이용해 DSV가 global explanation으로 작용할 수 있음을 보인다.}

% \subsection{Adversarial Attack}
% \jh{이것도 짧게 쓸수도, 아니면 안 쓸수도.. 렐웍에 넣을수도..}
% \subsection{Generative model}
% \jh{우리가 학습하는 것은 모델이 아니라 그것의 image이다. 따라서 우리의 모델은 generation model 중 하나로 볼 수 있는데, 기존의 generation model과 다른 점은 크게 두 가지로 요약할 수 있다. 첫 번째로 우리의 모델은 우리가 generate 하고자 하는 data manifold를 구할 수 없다.\hh{조금더 설명 추가가능? 많은 데이터셋( 데이터 매니폴드를 추정할 수 있는 어떤 무언가)가 없이도 generate 할 수 있다는 것인가?} 기존에 generative model로서 생각되는 gan, diffusion 모두 엄청나게 많은 양의 dataset을 전제한다. 또한, 일반적으로 genrative model이 구현하고자 하는 것은 true dataset distribution 그 자체이다. 하지만, 우리가 구현하고 싶은 것은 support vector로, true dataset distribution이 아닌, $\partial D$중, 즉 dataset의 boundary 중 가장 decision boundary를 잘 generate 하는것을 목적으로 하고 있다.}

